\par\medskip
\noindent
3. \underline{El caso de reducción semiestable. El teorema de reducción
semiestable.}

Damos primero un resultado auxiliar al estilo de SGA 3.

\medskip
\noindent
\underline{Proposición} 3.1. \underline{Sea $S$ un esquema, sea $G$ un
$S$-esquema en grupos conmutativo, plano, separado y de presentación
finita sobre $S$, sea $H$ un esquema en grupos cuasiseparado localmente
de tipo finito sobre $S$, sea $u:G\longrightarrow H$ un homomorfismo de
$S$-esquemas en grupos, y pongamos $K=\ker u$. Para todo $s\in S$, sea
$\bar s$ el espectro de una clausura algebraica de $k(s)$, y denotemos
por $\mu(s)$ (respectivamente, $\alpha(s)$) el rango reductivo
(respectivamente, el rango abeliano) de $K_s$, es decir, la dimensión
del mayor subtoro (respectivamente, del mayor cociente variedad
abeliana) de $(K_{\bar s})^o_{\mathrm{red}}$. Sean $s\in S$, $t\in S$
una generalización de $s$, y $p$ el exponente característico de $k(s)$.}

\smallskip
\noindent
a) \underline{Para todo entero $n>0$ primo con $p$, ${}_nK$ es étale,
de tipo finito y separado sobre $S$ en un entorno de $s$, y se tiene la
relación de divisibilidad}
\begin{equation}\tag{3.1.1}
 \operatorname{rank}({}_nK_s)\mid \operatorname{rank}({}_nK_t).
\end{equation}

\noindent
b) \underline{Se tiene}
\begin{equation}\tag{3.1.2}
 \mu(s)+2\alpha(s)\leq \mu(t)+2\alpha(t).
\end{equation}

\noindent
c) \underline{Si $(K_{\bar t})^o_{\mathrm{red}}$ es unipotente,
entonces $(K_{\bar s})^o_{\mathrm{red}}$ también lo es.}

\smallskip
\noindent
d) \underline{Supongamos que $G_s$ posee rango unipotente cero, que
$H$ es liso sobre $S$ y que $u_t$ es una isogenia. Entonces $u$ es
plano y cuasifinito sobre un entorno abierto $U$ de $s$; con mayor
razón, $K|_U$ es plano y cuasifinito.}

\smallskip
\noindent
e) \underline{Bajo las hipótesis de d), supongamos además que $u_t$ es
un monomorfismo (y, por tanto, un isomorfismo de $G_t$ sobre un subgrupo
abierto de $H_t$). Entonces $u_s$ es un isomorfismo de $G_s$ sobre un
subgrupo abierto de $H_s$. Si, además, $S$ es irreducible con punto
genérico $t$, existe un entorno abierto $U$ de $s$ tal que $u$ induce
un isomorfismo de $G|_U$ sobre un subgrupo abierto de $H|_U$.}

\medskip
\noindent
\underline{Demostración.}

\smallskip
\noindent
a) Podemos suponer que $n$ es invertible sobre $S$. Entonces
$n\,\mathrm{id}_G$ es étale sobre toda fibra (SGA 3, XV, 1.3); puesto
que $G$ es plano y de presentación finita sobre $S$, el criterio
fibrado de etalicidad (EGA IV, 17.8.2) muestra que
$n\,\mathrm{id}_G$ es étale. Por tanto, su núcleo ${}_nG$ es étale;
es también separado y de presentación finita sobre $S$, porque $G$ lo
es. Análogamente, ${}_nH$ es localmente de tipo finito, cuasiseparado y
formalmente no ramificado sobre $S$,\footnote{Al menos cuando $H$ es
conmutativo. En el caso general, ${}_nH$ sigue siendo no ramificado
sobre $S$ a lo largo de su sección unidad, que es todo lo que aquí se
necesita.} de modo que su sección unidad es una inmersión abierta
(EGA IV, 17.4.2 b')). En consecuencia, el núcleo de
\[
 {}_nu:{}_nG\longrightarrow {}_nH
\]
es un subesquema abierto de ${}_nG$, luego es étale sobre $S$; es
cuasicompacto en ${}_nG$ y, por tanto, de presentación finita sobre
$S$, y es separado sobre $S$ porque $G$ lo es. Este núcleo es
manifiestamente ${}_nK$.

La fórmula (3.1.1) se deduce de estas propiedades de ${}_nK$. En efecto,
tras un cambio de base podemos reducirnos al caso en que $S$ es local
henseliano con punto cerrado $s$, y utilizar entonces la descomposición
de ${}_nK$ en ${}_nK^f$ y un esquema de fibra especial vacía
(cf. 2.2.3). Más precisamente, el mismo argumento da un monomorfismo
sobre las fibras geométricas
\[
 {}_nK(\bar s)\lhook\joinrel\longrightarrow {}_nK(\bar t).
\]

\smallskip
\noindent
b) Elijamos un entero $n>0$ primo con el exponente característico de
$k(s)$ y con los rangos de los cocientes de las fibras geométricas de
$K$ en $s$ y $t$ por sus componentes neutras reducidas. Los dos miembros
de (3.1.1) se obtienen entonces elevando $n$, respectivamente, a
potencias cuyos exponentes son los dos miembros de (3.1.2), como se
deduce de la estructura conocida de los grupos algebraicos conmutativos.
Así, b) se deduce de a).

\smallskip
\noindent
c) La hipótesis de que $(K_{\bar t})^o_{\mathrm{red}}$ sea unipotente
equivale a las igualdades $\mu(t)=\alpha(t)=0$, y análogamente para
$K_{\bar s}$. La afirmación c) se deduce, pues, formalmente de b).

\smallskip
\noindent
d) Puesto que $u_t$ es una isogenia, $K_t$ es finito. Estamos, por
tanto, bajo las hipótesis de c), y $(K_{\bar s})^o_{\mathrm{red}}$ es
unipotente. Como $G_s$ posee rango unipotente cero, se sigue que $K_s$
es finito. Además, el hecho de que $u_t$ sea una isogenia da
$\dim G_t=\dim H_t$. Puesto que $G$ y $H$ son planos y localmente de
presentación finita, se tiene
\[
 \dim G_s=\dim G_t,\qquad \dim H_s=\dim H_t
 \quad\text{(SGA 3, VI$_B$, 4.3)},
\]
y, por tanto, $\dim G_s=\dim H_s$. Como $\dim\ker u_s=0$, se sigue que
$\dim H_s=\dim u_s(G_s)$ (SGA 3, VI$_B$, 1.2), lo cual, dado que $H_s$
es liso, equivale a que $u_s$ sea plano (SGA 3, VI$_A$, 5.6). Existe,
por tanto, un entorno abierto $U$ de $s$ tal que $u$ es plano y
cuasifinito en los puntos de la sección unidad sobre $U$ (SGA 3,
VI$_B$, 2.2(i) y 2.5(ii)). Como para $u_s$, se sigue que $u_{s'}$ es
plano para todo $s'\in U$; por consiguiente, $u|_{G_U}$ es plano y
cuasifinito, puesto que posee estas propiedades sobre toda fibra.

\smallskip
\noindent
e) Para la primera afirmación, un cambio de base nos reduce de inmediato
al caso en que $S$ es entero con punto genérico $t$, y por tanto a
demostrar la segunda afirmación de e). Esta se deduce inmediatamente
de d) y del lema siguiente.

\medskip
\noindent
\underline{Lema} 3.1.3. \underline{Sea $S$ un esquema, sea $K$ un
$S$-esquema cuasifinito de presentación finita, separado y plano sobre
$S$, sean $s\in S$ y $t\in\sisi{T}{S}$ una generalización de $s$. Entonces}
\[
 \operatorname{rank}K_s\leq \operatorname{rank}K_t.
\]
\underline{Si $\operatorname{rank}$} $K_\eta\leq1$ \underline{para todo punto máximo
$\eta$ de} $S$, \underline{entonces $K\longrightarrow S$ es una inmersión abierta
(y, por tanto, un isomorfismo si admite una sección).}

Para la primera afirmación, la localización estricta nos reduce al caso
en que $S$ es estrictamente local. Descomponiendo $K$ como unión
disjunta de un esquema finito $K^f$ y un esquema de fibra especial
vacía, obtenemos
\[
 \operatorname{rank}K_s
 =\operatorname{rank}K^f_s
 =\operatorname{rank}K^f_t
 \leq\operatorname{rank}K_t.
\]
La segunda igualdad se deduce de que $K^f$ es plano, finito y de
presentación finita sobre $S$ y, por tanto, localmente libre sobre $S$.

Para la segunda afirmación, el argumento precedente muestra ya que
$\operatorname{rank}K_s\leq1$ para todo $s\in S$. Así, $K_s$ es vacío
o isomorfo a $\operatorname{Spec}k(s)$. En consecuencia,
$K\longrightarrow S$ es un monomorfismo (SGA 3, VI$_B$, 2.11) y, al
ser plano y de presentación finita, es una inmersión abierta
(EGA IV, 17.9.1). Esto demuestra el lema.

\medskip
\noindent
\underline{Observación} 3.1.4. La parte 3.1 a), y el resto de la
Proposición 3.1, no proporcionan información alguna sobre ${}_nK_s$
cuando $n$ es una potencia de $p$ (para $p>1$). Al menos cuando $G$ es
liso (y $S$ es localmente noetheriano o $H$ es de presentación finita
sobre $S$), esencialmente el mismo método que en la demostración de a)
proporciona tal información. En efecto, reduciéndonos al caso en que
$S$ es local noetheriano completo y utilizando SGA 3, XV, 1.6 c), X,
3.2, IX, 3.6, y IX, 5.2, se obtiene
\begin{equation}\tag{3.1.4.1}
 \operatorname{rank}_{\mathrm{mult}}({}_nK_s)
 \mid
 \operatorname{rank}_{\mathrm{mult}}({}_nK_t),
\end{equation}
donde $\operatorname{rank}_{\mathrm{mult}}$ denota el rango del mayor
subgrupo de tipo multiplicativo del grupo algebraico conmutativo
considerado (SGA 3, XVII, 7.2.1). Obsérvese que un grupo algebraico de
torsión de tipo multiplicativo es finito. La fórmula (3.1.4.1) es
válida para todo entero $n\geq1$. Se deduce de SGA 3, XVII, 4.6.1(vi)
que, en c), si $K_t$ es unipotente (SGA 3, VII, 1.1), entonces $K_s$
también lo es.
